\scanpage{116}
% L1–2; large diagonal strokes through the entire draft
\deleted{Soit $X$ un espace topologique, \deleted{\unsure{dont chaque pt a un syst.|dont chaque pt $x$ a un syst.}}\inserted{ayant une base d'ouverts} \deleted{fondamental de voisinages} $U$ ayant les propriétés suivantes :
\margin{\unsure{Hyp|1\textsuperscript{ère}}}

% L3
(i) $U$ rétrocompact dans $X$, et $U$ quasi-compact

% L4–6
(ii) T\textsuperscript{te} \deleted{suite} famille $(E_i)_{i\in I}$ de parties \deleted{\unsure{fermées}} constructibles de $U$, recouvrant $U$, il y a une sous-famille finie recouvrant $U$.

% L7–11
Pour exprimer (ii), Introduisons [sur l'espace topologique $U$] la topologie dont les ouverts sont les \inserted{réunions quelconques de} parties loc. constructibles [c'est bien stable par intersections finies et réunions \unsure{qlc|qlq|quelc.}] donc une base d'ouverts formée des parties constructibles.
% L12–17
Lorsque $U$ est quasi-compact, et \deleted{que} \deleted{\unsure{$U$ qui admet|$U$ est admet}} admet une base d'ouverts rétrocompacts, ``loc constructible'' est d'ailleurs pareil que ``constructible'' ; Moyennant (i), les $U$ intervenant dans (Hyp) tombent dans le cas précédent.
\margin{$\varnothing$.}
% L17–19
Notons \inserted{que} si de plus $\deleted{\unsure{X}}U$ satisfait l'axiome $T_0$, \inserted{[c'est le cas des préschémas \unsure{\ldots}]} \deleted{\unsure{$U$ est}} alors une telle topologie est \uline{séparée}, et \unsure{même|en} totalement discontinue [pour deux pts $x\ne y$, il existe des ouverts $V\ni x$, $W\ni y$, $V\cap W=\varnothing$, $V\cup W=U$]
% L20–22
L'hyp (ii) exprime qu'elle est \uline{quasi-compacte} [donc \uline{compacte} dans le cas particulier que nous envisageons].}
% archive sheet number: 115
